% !PRINTED 25
\weberpage{25}
%<<CONT
from functions of one variable $x$ when the coefficients
$a_0,a_1\ldots a_n$ occurring in them are themselves regarded as
whole rational functions of other variables $y,z\ldots$. Thus
functions of $m$ variables arise from functions of $m-1$ variables.
All terms of such a function are of the form $x^r y^s z^t\cdots$,
multiplied by a coefficient; and if the function is properly
collected, then each combination of the exponents $r,s,t\ldots$
occurs only once. Two whole functions so ordered are regarded as equal
only when they contain the same products $x^r y^s z^t\cdots$ with the
same coefficients.

\begin{center}
{\large \S.\ 2.}

\vspace{0.5em}
{\large\gesperrt{A Theorem of Gauss.}}
\end{center}

We shall at once apply the multiplication rule for two whole rational
functions in order to prove a theorem of Gauss which will be useful to
us later, but which here is to serve as an introduction to the mode of
calculation and as an example.\footnote{Gauss, \emph{Disquisitiones
arithmeticae}, Art.~42.}

We consider here the case in which the coefficients in the functions
$A(x),B(x)$ are whole numbers, so that, by (4), the coefficients of
$C(x)=A(x)B(x)$ are also whole numbers.

If all integral coefficients $a_0,a_1,\ldots,a_m$ of a function
$A(x)$ have no common divisor, then the function $A(x)$ is called an
\gesperrt{original or primitive} function, and the theorem we wish to
prove is:

\gesperrt{If $A(x)$ and $B(x)$ are primitive functions, then their
product $C(x)$ is also a primitive function.}

The proof follows almost immediately from looking at formula (4).

For if all coefficients $c_0,c_1,c_2\ldots c_{m+n}$, as given in
(4), had a common divisor greater than $1$, then there would also have
to be at least one prime number which divides all these coefficients.
Let $p$ be such a prime number. By the assumption that $A(x),B(x)$ are
primitive, this prime can divide neither all $a_0,a_1\ldots a_m$ nor
all $b_0,b_1\ldots b_n$.
